\section{Related Work}
\label{sec:related_work}

\paragraph{Prompt auto tuning and agent-configuration search.}
GEPA demonstrates that trajectory feedback can guide reflective prompt evolution and outperform reinforcement learning on several language-agent tasks \citep{agrawal2025gepa}. ABSTRAL and EvoTest extend this idea from single prompts to multi-agent design documents and test-time agentic system evolution without gradients or fine-tuning \citep{song2026abstral,he2025evotest}. By treating language artifacts as optimizable objects, these methods can directly exploit execution feedback, but they mainly target prompts, system designs, or full configurations rather than reusable domain adaptation. \ourmethod{} instead optimizes a persistent skill document that can be trained, validated, exported, and reused with the adapted model, applying language-level controllability to a stable procedural skill state.

% \paragraph{Skill construction and skill evolution.}
\paragraph{Skill construction and skill evolution.}
SkillsBench and the SoK on agentic skills frame skills as reusable procedural knowledge, covering tool policies, applicability conditions, execution routines, and supporting resources \citep{li2026skillsbench,jiang2026sok}. Prior systems construct such skills from lifelong experience, trajectory lessons, skill knowledge bases, or heterogeneous domain resources \citep{yang2026autoskill,ni2026trace2skill,wang2026skillx,shen2026skillfoundry,memp}, and further refine them through failure analysis, creation-evaluation-revision loops, co-evolving generators and verifiers, collective updates, or reinforcement learning \citep{alzubi2026evoskill,luo2026skillforge,zhang2026coevoskills,ma2026skillclaw,xia2026skillrl,wang2025reinforcement,autorefine,procmem,evolver}. While these works emphasize skill discovery, repository growth, sharing, evolutionary search, or policy optimization, \ourmethod{} studies a narrower problem: how to train one compact domain skill with deep-learning-style controls such as trajectory batches, reflection minibatches, textual learning rates, validation gates, rejected-edit buffers, and slow/meta updates. This yields a controlled and auditable procedure for producing a portable \texttt{best\_skill.md} without changing model weights.
